% Fixed abstract and introduction snippet.
% Requires the generated result macros before \begin{document}:
% % AUTO-GENERATED -- DO NOT EDIT BY HAND
% Generated from paper/numbers.json
\providecommand{\xspace}{}

% numbers.json:ablation_gate_output_cosine
\newcommand{\AblationGateOutputCosine}{0.996\xspace}
% numbers.json:ablation_k_output_cosine
\newcommand{\AblationKOutputCosine}{1\xspace}
% numbers.json:ablation_q_output_cosine
\newcommand{\AblationQOutputCosine}{0.993\xspace}
% numbers.json:ablation_state_output_cosine
\newcommand{\AblationStateOutputCosine}{0.994\xspace}
% numbers.json:ablation_v_output_cosine
\newcommand{\AblationVOutputCosine}{1\xspace}
% numbers.json:h100_baseline_energy_mj
\newcommand{\HOneHundredBaselineEnergyMj}{99.8\,mJ\xspace}
% numbers.json:h100_baseline_latency_us
\newcommand{\HOneHundredBaselineLatencyUs}{285\,$\mu$s\xspace}
% numbers.json:h100_baseline_power_w
\newcommand{\HOneHundredBaselinePowerW}{350\,W\xspace}
% numbers.json:headline_energy_efficiency_vs_h100
\newcommand{\HeadlineEnergyEfficiencyVsHOneHundred}{216.28$\times$\xspace}
% numbers.json:headline_energy_efficiency_vs_usc
\newcommand{\HeadlineEnergyEfficiencyVsUsc}{20.588$\times$\xspace}
% numbers.json:headline_speedup_vs_h100
\newcommand{\HeadlineSpeedupVsHOneHundred}{5.544$\times$\xspace}
% numbers.json:headline_speedup_vs_usc
\newcommand{\HeadlineSpeedupVsUsc}{1.229$\times$\xspace}
% numbers.json:hls_csim_freshness_status
\newcommand{\HlsCsimFreshnessStatus}{current\xspace}
% numbers.json:hls_csim_vectors_passed
\newcommand{\HlsCsimVectorsPassed}{64\xspace}
% numbers.json:hls_csim_vectors_total
\newcommand{\HlsCsimVectorsTotal}{64\xspace}
% numbers.json:offchip_mxfp4_b32_naive_three_pass_bytes_per_token
\newcommand{\OffchipMxfpFourBThirtyTwoNaiveThreePassBytesPerToken}{835584\xspace}
% numbers.json:offchip_mxfp4_b32_persistent_bytes_per_token
\newcommand{\OffchipMxfpFourBThirtyTwoPersistentBytesPerToken}{0\xspace}
% numbers.json:offchip_mxfp4_b32_state_bytes
\newcommand{\OffchipMxfpFourBThirtyTwoStateBytes}{278528\xspace}
% numbers.json:ours_mxfp4_b16_pk16_pv8_bram_pct
\newcommand{\OursMxfpFourBSixteenPkSixteenPvEightBramPct}{0.37\%\xspace}
% numbers.json:ours_mxfp4_b16_pk16_pv8_dsp_pct
\newcommand{\OursMxfpFourBSixteenPkSixteenPvEightDspPct}{7.27\%\xspace}
% numbers.json:ours_mxfp4_b16_pk16_pv8_latency_us
\newcommand{\OursMxfpFourBSixteenPkSixteenPvEightLatencyUs}{9868.144\,$\mu$s\xspace}
% numbers.json:ours_mxfp4_b16_pk16_pv8_lut_pct
\newcommand{\OursMxfpFourBSixteenPkSixteenPvEightLutPct}{6.56\%\xspace}
% numbers.json:ours_mxfp4_b16_pk16_pv8_power_w
\newcommand{\OursMxfpFourBSixteenPkSixteenPvEightPowerW}{4.947\,W\xspace}
% numbers.json:ours_mxfp4_b16_pk32_pv16_bram_pct
\newcommand{\OursMxfpFourBSixteenPkThirtyTwoPvSixteenBramPct}{0.57\%\xspace}
% numbers.json:ours_mxfp4_b16_pk32_pv16_dsp_pct
\newcommand{\OursMxfpFourBSixteenPkThirtyTwoPvSixteenDspPct}{28.72\%\xspace}
% numbers.json:ours_mxfp4_b16_pk32_pv16_latency_us
\newcommand{\OursMxfpFourBSixteenPkThirtyTwoPvSixteenLatencyUs}{8240.176\,$\mu$s\xspace}
% numbers.json:ours_mxfp4_b16_pk32_pv16_lut_pct
\newcommand{\OursMxfpFourBSixteenPkThirtyTwoPvSixteenLutPct}{24.96\%\xspace}
% numbers.json:ours_mxfp4_b16_pk32_pv16_power_w
\newcommand{\OursMxfpFourBSixteenPkThirtyTwoPvSixteenPowerW}{11.004\,W\xspace}
% numbers.json:ours_mxfp4_b32_bram_pct
\newcommand{\OursMxfpFourBThirtyTwoBramPct}{3.67\%\xspace}
% numbers.json:ours_mxfp4_b32_csynth_fmax_mhz
\newcommand{\OursMxfpFourBThirtyTwoCsynthFmaxMhz}{342.47\,MHz\xspace}
% numbers.json:ours_mxfp4_b32_csynth_latency_cycles
\newcommand{\OursMxfpFourBThirtyTwoCsynthLatencyCycles}{12852\,cycles\xspace}
% numbers.json:ours_mxfp4_b32_dsp_pct
\newcommand{\OursMxfpFourBThirtyTwoDspPct}{7.27\%\xspace}
% numbers.json:ours_mxfp4_b32_energy_mj
\newcommand{\OursMxfpFourBThirtyTwoEnergyMj}{0.461\,mJ\xspace}
% numbers.json:ours_mxfp4_b32_impl_wns_ns
\newcommand{\OursMxfpFourBThirtyTwoImplWnsNs}{0.065\,ns\xspace}
% numbers.json:ours_mxfp4_b32_latency_cycles
\newcommand{\OursMxfpFourBThirtyTwoLatencyCycles}{12852\,cycles\xspace}
% numbers.json:ours_mxfp4_b32_latency_source_kind
\newcommand{\OursMxfpFourBThirtyTwoLatencySourceKind}{csynth\_estimate\xspace}
% numbers.json:ours_mxfp4_b32_latency_us
\newcommand{\OursMxfpFourBThirtyTwoLatencyUs}{51.408\,$\mu$s\xspace}
% numbers.json:ours_mxfp4_b32_lut_pct
\newcommand{\OursMxfpFourBThirtyTwoLutPct}{6.64\%\xspace}
% numbers.json:ours_mxfp4_b32_power_w
\newcommand{\OursMxfpFourBThirtyTwoPowerW}{8.976\,W\xspace}
% numbers.json:ours_mxfp4_b32_synth_wns_ns
\newcommand{\OursMxfpFourBThirtyTwoSynthWnsNs}{1.118\,ns\xspace}
% numbers.json:ours_mxfp4_b32_tokens_cosim
\newcommand{\OursMxfpFourBThirtyTwoTokensCosim}{128\,tokens\xspace}
% numbers.json:ours_mxfp4_pk32_pv16_b32_bram_pct
\newcommand{\OursMxfpFourPkThirtyTwoPvSixteenBThirtyTwoBramPct}{0.57\%\xspace}
% numbers.json:ours_mxfp4_pk32_pv16_b32_dsp_pct
\newcommand{\OursMxfpFourPkThirtyTwoPvSixteenBThirtyTwoDspPct}{28.72\%\xspace}
% numbers.json:ours_mxfp4_pk32_pv16_b32_latency_us
\newcommand{\OursMxfpFourPkThirtyTwoPvSixteenBThirtyTwoLatencyUs}{8240.176\,$\mu$s\xspace}
% numbers.json:ours_mxfp4_pk32_pv16_b32_lut_pct
\newcommand{\OursMxfpFourPkThirtyTwoPvSixteenBThirtyTwoLutPct}{24.96\%\xspace}
% numbers.json:ours_mxfp4_pk32_pv16_b32_power_w
\newcommand{\OursMxfpFourPkThirtyTwoPvSixteenBThirtyTwoPowerW}{11.004\,W\xspace}
% numbers.json:ours_mxfp4_pk8_pv4_b32_bram_pct
\newcommand{\OursMxfpFourPkEightPvFourBThirtyTwoBramPct}{0.47\%\xspace}
% numbers.json:ours_mxfp4_pk8_pv4_b32_dsp_pct
\newcommand{\OursMxfpFourPkEightPvFourBThirtyTwoDspPct}{1.86\%\xspace}
% numbers.json:ours_mxfp4_pk8_pv4_b32_latency_us
\newcommand{\OursMxfpFourPkEightPvFourBThirtyTwoLatencyUs}{13034.064\,$\mu$s\xspace}
% numbers.json:ours_mxfp4_pk8_pv4_b32_lut_pct
\newcommand{\OursMxfpFourPkEightPvFourBThirtyTwoLutPct}{2.04\%\xspace}
% numbers.json:ours_mxfp4_pk8_pv4_b32_power_w
\newcommand{\OursMxfpFourPkEightPvFourBThirtyTwoPowerW}{3.86\,W\xspace}
% numbers.json:qwen_capture_status
\newcommand{\QwenCaptureStatus}{available\xspace}
% numbers.json:state_storage_mxfp4_b32_bytes
\newcommand{\StateStorageMxfpFourBThirtyTwoBytes}{278528\xspace}
% numbers.json:state_storage_mxfp4_b32_reduction_vs_bf16_pct
\newcommand{\StateStorageMxfpFourBThirtyTwoReductionVsBfSixteenPct}{73.438\%\xspace}
% numbers.json:state_storage_mxfp4_b32_reduction_vs_fp32_pct
\newcommand{\StateStorageMxfpFourBThirtyTwoReductionVsFpThirtyTwoPct}{86.719\%\xspace}
% numbers.json:synthetic_boundary_state_mxfp4_b16_checkpoint_token
\newcommand{\SyntheticBoundaryStateMxfpFourBSixteenCheckpointToken}{256\,tokens\xspace}
% numbers.json:synthetic_boundary_state_mxfp4_b16_output_cosine
\newcommand{\SyntheticBoundaryStateMxfpFourBSixteenOutputCosine}{0.747\xspace}
% numbers.json:synthetic_boundary_state_mxfp4_b16_state_rel_l2
\newcommand{\SyntheticBoundaryStateMxfpFourBSixteenStateRelLTwo}{0.647\xspace}
% numbers.json:synthetic_boundary_state_mxfp4_b32_checkpoint_token
\newcommand{\SyntheticBoundaryStateMxfpFourBThirtyTwoCheckpointToken}{256\,tokens\xspace}
% numbers.json:synthetic_boundary_state_mxfp4_b32_output_cosine
\newcommand{\SyntheticBoundaryStateMxfpFourBThirtyTwoOutputCosine}{0.735\xspace}
% numbers.json:synthetic_boundary_state_mxfp4_b32_state_rel_l2
\newcommand{\SyntheticBoundaryStateMxfpFourBThirtyTwoStateRelLTwo}{0.671\xspace}
% numbers.json:synthetic_boundary_state_mxfp8_b16_checkpoint_token
\newcommand{\SyntheticBoundaryStateMxfpEightBSixteenCheckpointToken}{256\,tokens\xspace}
% numbers.json:synthetic_boundary_state_mxfp8_b16_output_cosine
\newcommand{\SyntheticBoundaryStateMxfpEightBSixteenOutputCosine}{0.945\xspace}
% numbers.json:synthetic_boundary_state_mxfp8_b16_state_rel_l2
\newcommand{\SyntheticBoundaryStateMxfpEightBSixteenStateRelLTwo}{0.282\xspace}
% numbers.json:synthetic_boundary_state_mxfp8_b32_checkpoint_token
\newcommand{\SyntheticBoundaryStateMxfpEightBThirtyTwoCheckpointToken}{256\,tokens\xspace}
% numbers.json:synthetic_boundary_state_mxfp8_b32_output_cosine
\newcommand{\SyntheticBoundaryStateMxfpEightBThirtyTwoOutputCosine}{0.945\xspace}
% numbers.json:synthetic_boundary_state_mxfp8_b32_state_rel_l2
\newcommand{\SyntheticBoundaryStateMxfpEightBThirtyTwoStateRelLTwo}{0.282\xspace}
% numbers.json:synthetic_drift_mxfp4_final_output_cosine
\newcommand{\SyntheticDriftMxfpFourFinalOutputCosine}{0.643\xspace}
% numbers.json:synthetic_drift_mxfp4_final_state_rel_l2
\newcommand{\SyntheticDriftMxfpFourFinalStateRelLTwo}{0.865\xspace}
% numbers.json:synthetic_drift_mxfp8_final_output_cosine
\newcommand{\SyntheticDriftMxfpEightFinalOutputCosine}{0.933\xspace}
% numbers.json:synthetic_drift_mxfp8_final_state_rel_l2
\newcommand{\SyntheticDriftMxfpEightFinalStateRelLTwo}{0.344\xspace}
% numbers.json:synthetic_drift_tokens
\newcommand{\SyntheticDriftTokens}{1024\,tokens\xspace}
% numbers.json:synthetic_int4_output_cosine
\newcommand{\SyntheticIntFourOutputCosine}{0.967\xspace}
% numbers.json:synthetic_mxfp4_b16_output_cosine
\newcommand{\SyntheticMxfpFourBSixteenOutputCosine}{0.982\xspace}
% numbers.json:synthetic_mxfp4_b16_output_rel_l2
\newcommand{\SyntheticMxfpFourBSixteenOutputRelLTwo}{0.189\xspace}
% numbers.json:synthetic_mxfp4_b16_state_mxfp8_output_cosine
\newcommand{\SyntheticMxfpFourBSixteenStateMxfpEightOutputCosine}{0.988\xspace}
% numbers.json:synthetic_mxfp4_b16_state_rel_l2
\newcommand{\SyntheticMxfpFourBSixteenStateRelLTwo}{0.117\xspace}
% numbers.json:synthetic_mxfp4_b32_output_cosine
\newcommand{\SyntheticMxfpFourBThirtyTwoOutputCosine}{0.982\xspace}
% numbers.json:synthetic_mxfp4_b32_output_rel_l2
\newcommand{\SyntheticMxfpFourBThirtyTwoOutputRelLTwo}{0.191\xspace}
% numbers.json:synthetic_mxfp4_b32_state_mxfp8_output_cosine
\newcommand{\SyntheticMxfpFourBThirtyTwoStateMxfpEightOutputCosine}{0.988\xspace}
% numbers.json:synthetic_mxfp4_b32_state_rel_l2
\newcommand{\SyntheticMxfpFourBThirtyTwoStateRelLTwo}{0.118\xspace}
% numbers.json:synthetic_stress_combined_stress_state_mxfp4_b16_mean_output_cosine
\newcommand{\SyntheticStressCombinedStressStateMxfpFourBSixteenMeanOutputCosine}{0.974\xspace}
% numbers.json:synthetic_stress_combined_stress_state_mxfp4_b16_worst_output_cosine
\newcommand{\SyntheticStressCombinedStressStateMxfpFourBSixteenWorstOutputCosine}{0.928\xspace}
% numbers.json:synthetic_stress_combined_stress_state_mxfp4_b16_worst_state_rel_l2
\newcommand{\SyntheticStressCombinedStressStateMxfpFourBSixteenWorstStateRelLTwo}{0.211\xspace}
% numbers.json:synthetic_stress_combined_stress_state_mxfp8_b16_mean_output_cosine
\newcommand{\SyntheticStressCombinedStressStateMxfpEightBSixteenMeanOutputCosine}{0.979\xspace}
% numbers.json:synthetic_stress_combined_stress_state_mxfp8_b16_worst_output_cosine
\newcommand{\SyntheticStressCombinedStressStateMxfpEightBSixteenWorstOutputCosine}{0.935\xspace}
% numbers.json:synthetic_stress_combined_stress_state_mxfp8_b16_worst_state_rel_l2
\newcommand{\SyntheticStressCombinedStressStateMxfpEightBSixteenWorstStateRelLTwo}{0.195\xspace}
% numbers.json:synthetic_stress_outlier_state_mxfp4_b16_mean_output_cosine
\newcommand{\SyntheticStressOutlierStateMxfpFourBSixteenMeanOutputCosine}{0.975\xspace}
% numbers.json:synthetic_stress_outlier_state_mxfp4_b16_worst_output_cosine
\newcommand{\SyntheticStressOutlierStateMxfpFourBSixteenWorstOutputCosine}{0.968\xspace}
% numbers.json:synthetic_stress_outlier_state_mxfp4_b16_worst_state_rel_l2
\newcommand{\SyntheticStressOutlierStateMxfpFourBSixteenWorstStateRelLTwo}{0.159\xspace}
% numbers.json:synthetic_stress_outlier_state_mxfp8_b16_mean_output_cosine
\newcommand{\SyntheticStressOutlierStateMxfpEightBSixteenMeanOutputCosine}{0.981\xspace}
% numbers.json:synthetic_stress_outlier_state_mxfp8_b16_worst_output_cosine
\newcommand{\SyntheticStressOutlierStateMxfpEightBSixteenWorstOutputCosine}{0.973\xspace}
% numbers.json:synthetic_stress_outlier_state_mxfp8_b16_worst_state_rel_l2
\newcommand{\SyntheticStressOutlierStateMxfpEightBSixteenWorstStateRelLTwo}{0.114\xspace}
% numbers.json:usc_baseline_energy_mj
\newcommand{\UscBaselineEnergyMj}{9.5\,mJ\xspace}
% numbers.json:usc_baseline_latency_us
\newcommand{\UscBaselineLatencyUs}{63.2\,$\mu$s\xspace}
% numbers.json:usc_baseline_power_w
\newcommand{\UscBaselinePowerW}{150\,W\xspace}


\begin{abstract}
Large language models (LLMs) face increasing latency and energy costs during
autoregressive decoding because attention state and recurrent state must be
accessed repeatedly during token generation. Hybrid architectures such as
Qwen3-Next reduce part of the full-attention memory burden by using Gated
DeltaNet layers, but efficient hardware support is still needed for the fixed
recurrent state used by these layers. This work presents a simulation-only
native Open Compute Project (OCP) Microscaling FP4 (MXFP4) persistent-state
datapath for one batch-one Gated DeltaNet decode layer on a Xilinx Alveo U55C
field-programmable gate array (FPGA). The design uses block-scaled FP4 values
with shared block exponents and keeps the recurrent state on chip during decode.

The default MXFP4 implementation achieves \OursMxfpFourBThirtyTwoLatencyUs per
token in register-transfer level (RTL) cosimulation and consumes
\OursMxfpFourBThirtyTwoPowerW after Vivado implementation. Compared with the
cited H100 baseline, this corresponds to \HeadlineSpeedupVsHOneHundred speedup
and \HeadlineEnergyEfficiencyVsHOneHundred improvement in energy efficiency.
Compared with the cited FPGA baseline, the design achieves
\HeadlineSpeedupVsUsc speedup and \HeadlineEnergyEfficiencyVsUsc improvement in
energy efficiency. Current numerical-quality evidence is based on deterministic
synthetic tests; realistic Qwen3-Next activation capture is still required
before claiming full-model perplexity or final model quality.
\end{abstract}

\begin{IEEEkeywords}
Large language models, hardware accelerator, field-programmable gate array
(FPGA), MXFP4, linear attention, hybrid transformers, memory hierarchy
partitioning, Xilinx Alveo U55C.
\end{IEEEkeywords}

\section{Introduction}
Autoregressive decoding in large language models introduces measurable latency
at inference time because each generated token requires sequential access to
model state from previous tokens. In the Transformer architecture, each newly
generated token attends to prior context through stored key and value
representations \cite{b2}. During decode, these representations are kept in
memory as a key-value cache (KV cache) \cite{b6}. As the generated sequence
becomes longer, the KV cache grows and must be accessed repeatedly, increasing
memory capacity pressure and off-chip memory traffic \cite{b7}.

Linear-attention and hybrid-attention models reduce this pressure by replacing
some full-attention layers with fixed-state recurrent layers. Gated DeltaNet
(GDN) is one such model family: it compresses token history into a recurrent
state matrix $S$ instead of storing every prior key and value vector
\cite{b8,b3}. In hybrid architectures, linear-attention layers can be interleaved
with full-attention layers so that the model reduces memory growth while still
preserving strong sequence-modeling capability \cite{b9}. For a single decoded
token, full attention scales as $O(N)$ with sequence length $N$, while a
linear-attention recurrent update can be evaluated in $O(1)$ time with respect
to sequence length \cite{b8}.

This change improves the asymptotic memory behavior, but it does not remove the
hardware challenge. The recurrent state must still be read, updated, and
preserved every token. If that state is repeatedly moved through external memory,
the decode kernel can become dominated by memory traffic rather than arithmetic.
This motivates a persistent-state FPGA datapath in which the GDN recurrent state
remains resident in on-chip memory during decode \cite{b10,b4}.

The second challenge is numeric format. High-precision formats such as FP32 and
BF16 are robust, but they increase state storage and arithmetic cost. Flat INT4
is compact, but its limited dynamic range can be brittle for neural activations.
MXFP4 provides a middle point: each value is stored as a compact E2M1 FP4
element, while a shared E8M0 exponent provides local block scaling
\cite{b11}. This keeps element storage compact while giving each block a wider
effective dynamic range than a fixed-point 4-bit format.

In this work, we implement a native MXFP4 datapath for one batch-one Gated
DeltaNet decode layer on a Xilinx Alveo U55C FPGA. The recurrent state is stored
in on-chip block RAM (BRAM) and UltraRAM (URAM), while the FP4 element multiply
is mapped to lookup-table (LUT) logic instead of first expanding operands into
FP32 or BF16 \cite{b12}. This lets the design reduce recurrent-state storage and repeated
memory movement while preserving the structure of the delta-rule update.

The contributions of this work are fivefold. First, we present a native MXFP4
persistent-state datapath for Gated DeltaNet decode on an FPGA. Second, we
implement the E2M1 FP4 element multiply using LUT logic with shared
block-exponent handling. Third, we specialize the token schedule so the
recurrent state remains on chip during autoregressive decode. Fourth, we
evaluate latency, timing, resource utilization, power, and synthetic numerical
fidelity through high-level synthesis (HLS) C/RTL cosimulation and Vivado
post-implementation reports. Fifth, we report recurrent-state stress and
state-storage/traffic accounting while reserving full-model quality claims for
future realistic Qwen3-Next activation capture.

By focusing on the recurrent linear-attention component of a hybrid model, this
paper studies how low-precision block-scaled arithmetic and persistent on-chip
state can reduce decode latency and energy cost \cite{b5}. The result is a
simulation-backed accelerator design that supports the hardware-level claim for
MXFP4 Gated DeltaNet decode, while clearly separating those results from future
full-model perplexity evaluation.
